%%%%%%%%%%%%%%%%%%%%%%%%%%%%%%%%%%%%%%%%%%%%%%%%%%%%%%%%%%%%%%%%%%%%%%%%%%%%%%%%
%%%%%%%%%%%%%%%%%%   Vorlage für eine Abschlussarbeit   %%%%%%%%%%%%%%%%%%%%%%%%
%%%%%%%%%%%%%%%%%%%%%%%%%%%%%%%%%%%%%%%%%%%%%%%%%%%%%%%%%%%%%%%%%%%%%%%%%%%%%%%%

% Erstellt von Maximilian Nöthe, <maximilian.noethe@tu-dortmund.de>
% ausgelegt für lualatex und Biblatex mit biber

% Kompilieren mit
% lualatex dateiname.tex
% biber dateiname.bcf
% lualatex dateiname.tex
% lualatex dateiname.tex
% oder einfach mit:
% make

\documentclass[
  tucolor,
  BCOR=12mm,     % 12mm binding corrections, adjust to fit your binding
  parskip=half,  % new paragraphs start with half line vertical space
  open=any,      % chapters start on both odd and even pages
  cleardoublepage=plain,  % no header/footer on blank pages
]{tudothesis}


% Warning, if another latex run is needed
\usepackage[aux]{rerunfilecheck}

% just list chapters and sections in the toc, not subsections or smaller
\setcounter{tocdepth}{1}

%------------------------------------------------------------------------------
%------------------------------ Sprache und Schrift: --------------------------
%------------------------------------------------------------------------------
\usepackage{fontspec}
\defaultfontfeatures{Ligatures=TeX}  % -- becomes en-dash etc.

% german language
\usepackage{polyglossia}
\setdefaultlanguage{german}

% for english abstract and english titles in the toc
\setotherlanguages{english}

% intelligent quotation marks, language and nesting sensitive
\usepackage[autostyle]{csquotes}

% microtypographical features, makes the text look nicer on the small scale
\usepackage{microtype}

%------------------------------------------------------------------------------
%------------------------ Für die Matheumgebung--------------------------------
%------------------------------------------------------------------------------

\usepackage{amsmath}
\usepackage{amssymb}
\usepackage{mathtools}
\usepackage{braket}

% Enable Unicode-Math and follow the ISO-Standards for typesetting math
\usepackage[
  math-style=ISO,
  bold-style=ISO,
  sans-style=italic,
  nabla=upright,
  partial=upright,
]{unicode-math}
\setmathfont{Latin Modern Math}

% nice, small fracs for the text with \sfrac{}{}
\usepackage{xfrac}


%------------------------------------------------------------------------------
%---------------------------- Numbers and Units -------------------------------
%------------------------------------------------------------------------------

\usepackage[
  locale=DE,
  separate-uncertainty=true,
  per-mode=symbol-or-fraction,
]{siunitx}
\sisetup{math-micro=\text{µ},text-micro=µ}

%------------------------------------------------------------------------------
%-------------------------------- tables  -------------------------------------
%------------------------------------------------------------------------------

\usepackage{booktabs}       % stellt \toprule, \midrule, \bottomrule

%------------------------------------------------------------------------------
%-------------------------------- graphics -------------------------------------
%------------------------------------------------------------------------------

\usepackage{graphicx}
\usepackage{grffile}

% allow figures to be placed in the running text by default:
\usepackage{scrhack}
\usepackage{float}
\floatplacement{figure}{htbp}
\floatplacement{table}{htbp}

% keep figures and tables in the section
\usepackage[section, below]{placeins}


%------------------------------------------------------------------------------
%---------------------- customize list environments ---------------------------
%------------------------------------------------------------------------------

\usepackage{enumitem}

%------------------------------------------------------------------------------
%------------------------------ Bibliographie ---------------------------------
%------------------------------------------------------------------------------

\usepackage[
  backend=biber,   % use modern biber backend
  autolang=hyphen, % load hyphenation rules for if language of bibentry is not
                   % german, has to be loaded with \setotherlanguages
                   % in the references.bib use langid={en} for english sources
]{biblatex}
\addbibresource{references.bib}  % die Bibliographie einbinden
\DefineBibliographyStrings{german}{andothers = {{et\,al\adddot}}}

%------------------------------------------------------------------------------
%------------------------------ Sonstiges: ------------------------------------
%------------------------------------------------------------------------------

\usepackage[pdfusetitle,unicode,linkbordercolor=tugreen]{hyperref}
\usepackage{bookmark}
\usepackage[shortcuts]{extdash}

%------------------------------------------------------------------------------
%-------------------------    Angaben zur Arbeit   ----------------------------
%------------------------------------------------------------------------------

\author{Maximilian Nöthe}
\title{\LaTeX-Dokumentenklasse und Vorlage für Abschlussarbeiten an der TU Dortmund}
\date{2014}
\birthplace{Castrop-Rauxel}
\chair{Lehrstuhl für Experimentelle Physik V}
\division{Fakultät Physik}
\thesisclass{Bachelor of Science}
\submissiondate{31. September 2015}
\firstcorrector{Prof.~Dr.~Erstgutachter}
\secondcorrector{Prof.~Dr.~Zweitgutachter}

% tu logo on top of the titlepage
\titlehead{\includegraphics[height=1.5cm]{logos/tu-logo.pdf}}

\begin{document}
\frontmatter
\thispagestyle{empty}
\setcounter{page}{2}
\section*{Hinweise}
Empfohlen wird die Verwendung dieser Vorlage mit der jeweils aktuellsten TeXLive Version (Linux, Windows) bzw. MacTeX Version (MacOS).
Aktuell ist dies TeXLive 2016. Download hier:
\begin{center}
  \ttfamily\url{https://www.tug.org/texlive/}
\end{center}
Bei Verwendung von TexLive Versionen 2014 und älter sollte
die Zeile
\begin{center}
\verb+\RequirePackage{fixltx2e}+ 
\end{center}
als erste Zeile der Präambel noch vor der Dokumentenklasse eingefügt werden.
Dies lädt diverse Bugfixes für LaTeX, die ab TexLive 2015 Standard sind.

Die Vorlage \texttt{thesis.tex} ist für die Kompilierung mit \texttt{lualatex} ausgelegt, mit wenigen Anpassungen kann sie aber auch mit \texttt{pdflatex} oder \texttt{xelatex} verwendet werden. 
Die Dokumentenklasse \texttt{tudothesis.cls} kann mit allen drei Programmen verwednet werden.

Achten Sie auch auf die Kodierung der Quelldateien.
Bei Verwendung von Xe\LaTeX\ oder Lua\LaTeX\ (empfohlen) müssen die
Quelldateien UTF-8 kodiert sein.
Bei Verwendung von pdf\LaTeX\ nutzen Sie die Pakete \texttt{inputenc} und \texttt{fontenc} mit der korrekten Wahl der Kodierungen.

Eine aktuelle Version dieser Vorlage steht unter 
\begin{center}
  \ttfamily\url{https://github.com/maxnoe/tudothesis}
\end{center}
zur Verfügung.

Alle verwendeten Pakete werden im \LaTeX{} Kurs von Pep et al.\ erklärt:
\begin{center}
  \ttfamily\url{http://toolbox.pep-dortmund.org/notes}
\end{center}

Für Rückmeldungen und bei Problemen mit der Klasse oder Vorlage, bitte ein \emph{Issue} auf GitHub aufmachen oder eine Email an
\href{mailto:maximilian.noethe@tu-dortmund.de}{maximilian.noethe@tu-dortmund.de} schreiben.

Wenn Sie die Dokumentenklasse mit der Option \texttt{tucolor} laden, werden verschiedene Elemente in TU-Grün gesetzt.

\maketitle

% Gutachterseite
\makecorrectorpage

% hier beginnt der Vorspann, nummeriert in römischen Zahlen
\thispagestyle{plain}

\section*{Kurzfassung}
Hier steht eine Kurzfassung der Arbeit in deutscher Sprache inklusive der Zusammenfassung der
Ergebnisse.
Zusammen mit der englischen Zusammenfassung muss sie auf diese Seite passen.

\section*{Abstract}
\begin{english}
The abstract is a short summary of the thesis in English, together with the German summary it has to fit on this page.
\end{english}

\tableofcontents

\mainmatter
% Hier beginnt der Inhalt mit Seite 1 in arabischen Ziffern
\chapter{Einleitung}
Hier folgt eine kurze Einleitung in die Thematik der Bachelorarbeit.
Die Einleitung muss kurz sein, damit die vorgegebene Gesamtlänge der 
Arbeit von 25 Seiten nicht überschritten wird. 
Die Beschränkung der Seitenzahl sollte man ernst nehmen,
da Überschreitung zu Abzügen in der Note führen kann. 
Um der Längenbeschränkung zu genügen, darf auch nicht an der Schriftgröße,
dem Zeilenabstand oder dem Satzspiegel (bedruckte Fläche der Seite) manipuliert werden.

\chapter{Theoretische Grundlagen}


\section{Floquet-Theorie}
  Die Floquet-Theorie ist ein nützliches Werkzeug zur der Lösung von quantenmechanischen Systemen, welche durch einen zeitlich periodischen Hamilton-Operator
  \begin{equation}
    H(t) = H(t+T) \; ,
  \end{equation}
  mit der Periode $T$, beschrieben werden.

  Das Floquet-Theorem besagt, dass bei einem solchen System, die Lösungen $\Psi_n(x,t)$ der Schrödinger-Gleichung
  \begin{equation}
    \text{i}\hbar\frac{\partial}{\partial t}\Psi_n(x,t) = H(t)\Psi_n(x,t)
    \label{schroedinger}
  \end{equation}
  in Ortsdarstellung die Form
  \begin{equation}
    \Psi_n(x,t) = \text{e}^{-\frac{i}{\hbar}\epsilon_nt}\Phi_n(x,t)
    \label{floquet_theorem}
  \end{equation}
  haben.
  Hierbei sind $\Phi_n(x,t) = \Phi_n(x,t+T)$ $T$-periodische Funktionen, die sogenannten Floquet-Moden, und $\epsilon_n$ die zugehörigen reellen Quasienergien, wobei diese Bezeichnungen gewählt wurden aufgrund der Parallele zu den Bloch-Moden und Quasiimpulsen des Bloch-Theorems \cite{haengi}.
  Das Floquet-Theorem kann damit als "Bloch-Theorem in der Zeit" aufgefasst werden \cite{sherly}.

  Durch Einsetzen dieses Ansatzes für die Wellenfunktionen (\ref{floquet_theorem}) in die Schrödingergleichung (\ref{schroedinger}) erhalten wir
  \begin{equation}
    \epsilon_n \Phi_n(x,t) = \left(H(t)-\text{i}\hbar\frac{\partial}{\partial t}\right)\Phi_n(x,t) = \cal{H}(t) \Phi_n(x,t) \; .
    \label{eigenwertproblem}
  \end{equation}
  Die Lösung der Schrödinger Gleichung konnte somit auf die Lösung eines Eigenwertproblems für den neuen Operator $\cal{H}(t)$ zurückgeführt werden \cite{sherly}.

  Die hermitischen Operatoren $H(t)$ und $\cal H(t)$ operieren auf dem Hilbertraum $\cal{L}^2 \otimes \cal{T}$.
  Dabei ist $\cal{L}^2$ der Raum der quadratintegrablen Funktionen und $\cal{T}$ der Raum der auf $[0,T]$ integrablen Funktionen, da die Operatoren $T$-periodisch sind \cite{haenggi}.
  Nach dem Spektralsatz bilden die Eigenfunktionen $\Phi_n(x,t)$ von $\cal{H}(t)$ eine Orthogonalbasis von $\cal{L}^2 \otimes \cal{T}$, welche auf eine Orthonormalbasis normiert werden kann, wodurch wir das Skalarprodukt definieren k"onnen als:
  %Es ergibt sich
  \begin{align}
    \begin{split}
    \braket{\braket{\Phi_n(x,t)|\Phi_m(x,t)}} &= \frac{1}{T} \int_0^T \braket{\Phi_n(x,t)|\Phi_n(x,t)} \text{d}t \\
    &= \frac{1}{T} \int_0^T \int_{-\infty}^{\infty} \Phi_n^*(x,t)\Phi_m(x,t) \: \text{d}x\text{d}t = \delta_{n,m} \; .
    \label{skalarprodukt_einzelner}
    \end{split}
  \end{align}



  \subsection{\texorpdfstring{Zeitlich gemittelter Erwartungswert der Energie $\bar{H}_n$}{Zeitlich gemittelter Erwartungswert der Energie bar{H}_n}}

    Da $H(t)$ nicht zeitlich konstant ist, sind auch dessen Erwartungswerte, die Energien des Systems, zeitabhängig.
    Ein direkter Vorteil der Floquet-Theorie liegt darin, dass sich die durchschnittliche Energie
    \begin{equation}
      \bar{H}_n  = \braket{\braket{\Psi_n(x,t)|H(t)|\Psi_m(x,t)}}
      %\label{mittleres_H}
    \end{equation}
    des $n$-ten Zustandes $\Psi_n(x,t)$ leicht über die Quasienergien $\epsilon_n$ berechnen lässt, ohne explizit Integrale zu lösen.

    Dazu ersetzen wir $H(t)$ mit Hilfe von (\ref{eigenwertproblem}).
    Außerdem unterscheiden sich die Floquet-Moden $\Phi_n(x,t)$ und die Wellenfunktionen $\Psi_n(x,t)$ nur durch eine komplexe Phase, daher sind deren Skalarprodukte identisch.
    Weiterhin benutzen wir (\ref{eigenwertproblem}), dass die Floquet-Moden Eigenfunktionen von $\cal{H}(t)$ sind:
    \begin{align}
      \begin{split}
      \bar{H}_n  &= \braket{\braket{\Phi_n(x,t)|\cal{H}(t)+\text{i}\hbar\frac{\partial}{\partial t}|\Phi_n(x,t)}} \\
      &=\epsilon_n \braket{\braket{\Phi_n(x,t)|\Phi_n(x,t)}} + \braket{\braket{\Phi_n(x,t)|\text{i}\hbar\frac{\partial}{\partial t}|\Phi_n(x,t)}} \\
      &= \epsilon_n + \braket{\braket{\Phi_n(x,t)|\text{i}\hbar\frac{\partial}{\partial t}|\Phi_n(x,t)}} \; .
    \end{split}
    \end{align}

    Nähere Betrachtung zeigt, dass
    \begin{equation}
      \text{i}\hbar\frac{\partial}{\partial t} = -\omega \frac{\partial \cal{H}(t)}{\partial \omega}
    \end{equation}
    gilt \cite{haenggi}.
    Da für $\epsilon_n$ und $\cal H(t)$ mit (\ref{eigenwertproblem}) eine Eigenwertgleichung vorliegt, kann das Hellman-Feynman-Theorem angewendet werden \cite{hellmann online quelle}.
    Dieses gibt eine Verbindung zwischen den Ableitungen der Eigenwerte und der Ableitung des Hamilton-Operators an.
    In unserem Fall erhalten wir dadurch
    \begin{equation}
      \frac{\partial \epsilon_n}{\partial \omega} = \braket{\braket{\Phi_n(x,t)|\frac{\partial \cal H(t)}{\partial \omega}|\Phi_n(x,t)}} \; .
    \end{equation}
    Damit folgt \cite{haenggi}:
    \begin{equation}
      \bar{H}_n = \epsilon_n - \omega\frac{\partial \epsilon_n}{\partial \omega} \; .
      \label{mittleres_H}
    \end{equation}

  \newpage



  \section{Allgemeine Lösung der Schrödinger-Gleichung des getriebenen harmonischen Oszillators in der Quantenmechanik}
    \label{lsg_einzelner}
    Der Hamilton-Operator eines harmonischen Oszillators der Masse $m$, welcher mit einer beliebigen aber periodischen äußeren Kraft $S(t)=S(t+T)$ getrieben wird, hat die Form
    \begin{equation}
      H(t) = H(t+T) = \frac{p^2}{2m} + \frac{1}{2}m\omega_0^2x^2-S(t)x \; ,
    \end{equation}
    mit $\omega_0=\sqrt{k/m}$.
    Dieses System kann exakt gelöst werden, indem die Schrödinger-Gleichung durch einen Variablenwechsel und zwei unitäre Transformationen auf die bekannte Form des ungetriebenen Oszillators reduziert wird \cite{haenggi}.

    \textbf{1) Variablenwechsel}\\
    Für den neuen Ortsoperator bzw. die Ortsvariable wird eine zeitabhängige Verschiebung angesetzt:
    \begin{equation}
      x \rightarrow y=x-\zeta(t) \; ;
    \end{equation}
    Wie zu erwarten verändert sich der Impuls(operator) durch die Translation im Ort nicht, da $\zeta(t)$ bei der Ortsableitung wegfällt.
    %wobei der Impuls(operator) unverändert bleibt.
    %Der Impulsoperator bleibt dadurch unverändert.

    Mit der neuen Zeitableitung der Wellenfunktion
    \begin{equation}
      \text{i}\hbar \frac{\partial}{\partial t} \Psi(y(t),t) = \text{i}\hbar \dot{\Psi} -\dot{\zeta}\frac{\partial}{\partial y}\Psi(y(t),t)
    \end{equation}
    wird die Schrödinger-Gleichung zu:
    \begin{equation}
      \text i \hbar \dot{\Psi}(y,t) = \left[\text i \hbar \dot{\zeta}\frac{\partial}{\partial y}-\frac{\hbar^2}{2m}\frac{\partial^2}{\partial y^2}+\frac{1}{2}m\omega_0^2(y+\zeta)^2-(y+\zeta)S(t)\right]\Psi(y,t) \; .
      \label{schroedinger_einzeln_getrieben}
    \end{equation}

    \textbf{2) Unitäre Tranformation für $\Psi(y,t)$}\\
    Im Weiteren wählen wir die unitäre Transformation
    \begin{equation}
      \Psi(y,t) = \text e^{\frac{\text i}{\hbar}m\dot \zeta y}\Lambda(y,t) \; .
    \end{equation}
    Durch Einsetzen in die Schrödinger-Gleichung (\ref{schroedinger_einzeln_getrieben}) und Ausrechnen der Ableitungen erhalten wir
    \begin{align}
      \begin{split}
        &\text e^{\frac{\text i}{\hbar}m\dot \zeta y}(\text i \hbar \dot \Lambda(y,t) - my \ddot{\zeta}\Lambda(y,t)) =\\
         &\text e^{\frac{\text i}{\hbar}m\dot \zeta y} \left[\left(-m\dot \zeta^2 \Lambda(y,t) + \text i \hbar \dot \zeta \frac{\partial}{\partial y} \Lambda(y,t) \right) \right. \\
         &\left. + \left(\frac{1}{2}m\dot \zeta^2 \Lambda(y,t) - \text{i} \hbar \dot \zeta \frac{\partial}{\partial y} \Lambda(y,t) - \frac{\hbar}{2m}\frac{\partial^2}{\partial y^2} \Lambda(y,t)  \right)
        + \frac{1}{2}m\omega_0^2(y+\zeta)^2  + (y+\zeta)S(t)  \right] \; .
      \end{split}
    \end{align}
    \iffalse
    Durch Einsetzen in die Schrödinger-Gleichung und Ausrechnen der Ableitungen erhalten wir für die linke Seite von (\ref{schroedinger_einzeln_getrieben})
    \text e^{\frac{\text i}{\hbar}m\dot \zeta y}(\text i \hbar \dot \Lambda(y,t) - my \ddot{\zeta}\Lambda(y,t))
    \begin{equation}
    \end{equation}
    und für die rechte Seite
    \begin{align}
      \begin{split}
      \text e^{\frac{\text i}{\hbar}m\dot \zeta y} \left[\left(-m\dot \zeta^2 \Lambda(y,t) + \text i \hbar \dot \zeta \frac{\partial}{\partial y} \Lambda(y,t) \right) + \left(\frac{1}{2}m\dot \zeta^2 \Lambda(y,t) - \text{i} \hbar \dot \zeta \frac{\partial}{\partial y} \Lambda(y,t) - \frac{\hbar}{2m}\frac{\partial^2}{\partial y^2} \Lambda(y,t)  \right) \right. \\
      \left. + \left(\frac{1}{2}m\omega_0^2y^2\Lambda(y,t) + m\omega_0^2y\zeta\Lambda(y,t) + \frac{1}{2}m\omega_0^2 \zeta^2\Lambda(y,t)) \right) + \left(-yS(t)\Lambda(y,t) - \zeta S(t)\Lambda(y,t) \right) \right]
    \end{split}
    \end{align}
    \fi

    \newpage

    Indem auf beiden Seiten durch die Exponentialfunktion geteilt wird, bekommen wir eine Differentialgleichung für $\Lambda(y,t)$.
    Außerdem können durch geschicktes Umsortieren der Terme die Lagrange-Funktion $L(\zeta,\dot \zeta, t)$
    \begin{equation}
      L(\zeta,\dot \zeta, t) = \frac{1}{2}m\dot \zeta^2 - \frac{1}{2}m\omega_0^2\zeta^2 + S(t)\zeta
      \label{lagrange_zeta}
    \end{equation}
      sowie die Bewegungsgleichung des klassischen getriebenen harmonischen  Oszillators \cite{husimi}
      \begin{equation}
        m\ddot \zeta + m\omega_0^2\zeta - S(t) = 0
        \label{dgl_zeta}
      \end{equation}
    für die Verschiebung $\zeta(t)$ identifiziert werden.
    Die Differentialgleichung für $\Lambda(y,t)$ ist folglich
    \begin{equation}
      \text i \hbar \dot \Lambda(y,t) = \left[ \frac{-\hbar^2}{2m}\frac{\partial^2}{\partial y^2} + \frac{1}{2}m\omega_0^2y^2 + (m\ddot \zeta + m\omega_0^2y\zeta - S(t))y - L(\zeta,\dot \zeta, t) \right]\Lambda(y,t) \; .
      \label{dgl_lambda}
    \end{equation}
    Um die Gleichung zu vereinfachen, wählen wir $\zeta(t)$ nun so, dass es gerade die klassische Bewegungsgleichung erfüllt, der entsprechende Term in (\ref{dgl_lambda}) also verschwindet.
    Nur noch die Lagrange-Funktion unterscheidet diese Differential-Gleichung von der des ungetriebenen Oszillators.

    \textbf{3) Unitäre Transformation für $\Lambda(y,t)$}\\
    Zuletzt wählen wir den Ansatz
    \begin{equation}
      \Lambda(y,t) = \text e^{\frac{\text i}{\hbar}\int_0^tL \: \text d t'}\chi(y,t)
    \end{equation}
    für $\Lambda(y,t)$, um die Lagrange-Funktion in (\ref{dgl_lambda}) zu eliminieren.
    Dadurch wird diese Differential-Gleichung für $\Lambda(y,t)$ bzw. die ursprüngliche Schrödinger-Gleichung auf einen ungetriebenen Oszillator für $\chi(y,t)$ zurückgeführt:
    \begin{equation}
      \text i \hbar \dot \chi(y,t) = \left[ \frac{-\hbar^2}{2m}\frac{\partial^2}{\partial y^2} + \frac{1}{2}m\omega_0^2y^2 \right]\chi(y,t) \; .
    \end{equation}
    Das bedeutet die $\chi_n(y,t)$ sind die Wellenfunktionen des ungetriebenen Oszillators und die Gesamtlösung der Schrödinger-Gleichung des getriebenen Oszillators ist damit gegeben durch
    \begin{align}
      \begin{split}
      \Psi_n(x,t) &= \Psi_n(y=x-\zeta(t),t) \\
      &= N_nO_n\left(\sqrt{\frac{m\omega_0}{\hbar}}(x-\zeta(t))\right) \text e^{\frac{-m\omega_0}{2\hbar}(x-\zeta(t))^2} \\
      &\quad \cdot \text e^{\frac{\text i}{\hbar}\left(m\dot \zeta(t)(x-\zeta(t))-E_nt+\int_0^tL(\dot \zeta,\zeta,t')\:\text dt'\right)} \; , \\
      &\quad \quad n \in \mathbb{N}_0 \; .
    \end{split}
    \end{align}
    Dabei sind $O_n$ die Hermit-Polynome, $E_n = \hbar \omega_0(n+1/2)$ die bekannten Eigenenergien der Zustände des ungetriebenen Oszillators und
    \begin{equation}
      N_n = \left(\frac{m\omega_0}{\pi \hbar}\right) \frac{1}{\sqrt{2^nn!}}
    \end{equation}
    dessen Normierungsfaktoren.
    Die Lösungen des getriebenen Oszillators sind somit bezüglich des Skalarproduktes (\ref{skalarprodukt_einzelner}) normiert.

    Die Lösung entspricht damit einem, um die klassiche Lösung $\zeta(t)$ verschobenen, ungetriebenen Oszillator, mit einer zusätzlichen zeit- und ortsabhängigen komplexen Phase.
    Die treibende Kraft $S(t)$ geht in die klassiche Lösung $\zeta(t)$ und direkt, über das Wirkungsintegral, in die komplexe Phase ein.


    \subsection{\texorpdfstring{Identifizierung der Quasienergien $\epsilon_n$ und Floquet-Moden $\Phi_n(x,t)$}{Identifizierung der Quasienergien \epsilon_n und Floquet-Moden \Phi_n(x,t)}}
      Nach dem Floquet-Theorem für periodische Hamilton-Operatoren (\ref{floquet_theorem}) kann die Lösung des getriebenen Oszillators geschrieben werden als
      \begin{equation}
        \Psi_n(x,t) = \text e^{\frac{-\text i}{\hbar}\epsilon_nt}\Phi_n(x,t) \; ,
      \end{equation}
      mit $\Phi_n(x,t)=\Phi_n(x,t+T)$.
      Nun können wir alle $T$-periodischen Terme als $\Phi_n(x,t)$ identifizieren, und alle Terme im Exponenten, die linear in $t$ sind, als $-\text i\epsilon_n t/\hbar$ \cite{haenggi}.

      Alle Funktionen von $(x-\zeta(t))$ haben die Periode $T$, da $\zeta(T)$ als Lösung der klassichen Bewegungsgleichung mit $S(t)=S(t+T)$ die Periode $T$ hat. Das
      Ergebnis des Integrals über die Lagrange-Funktion kann nur $T$-periodisch oder linear in $t$ sein, deshalb sind die Quasienergien gegeben durch die $E_n$ und den linearen Teil des Integrals:
      \begin{equation}
        \epsilon_n = E_n - \frac{1}{T} \int_0^T L(\dot \zeta, \zeta, t) \: \text d t \; .
      \end{equation}
      Die Floquet-Moden sind demnach
      \begin{align}
        \begin{split}
          \Phi_n(x,t) &=
           N_nH_n\left(\sqrt{\frac{m\omega_0}{\hbar}}(x-\zeta(t))\right) \text e^{\frac{-m\omega_0}{2\hbar}(x-\zeta(t))^2} \\
          &\quad \cdot \text e^{\frac{\text i}{\hbar}\left(m\dot \zeta(t)(x-\zeta(t))+\int_0^tL(\dot \zeta,\zeta,t')\:\text d t-\frac{t}{T} \int_0^T L(\dot \zeta,\zeta,t)\:\text dt\: \right)} \; , \\
          &\quad \quad n \in \mathbb{N}_0 \; .
        \end{split}
      \end{align}


    \subsection{Quasienergien für eine sinusoidiale Treibkraft}
      Hier wird ein Beispiel einer treibenden Kraft diskutiert \cite{haenggi}:
      \begin{equation}
        S(t) = S(t+T) = A\sin(\omega t),
      \end{equation}
      wobei $T=2\pi / \omega$ ist.
      Setzen wir die allgemeine homogene Lösung gleich null, wird die Lösung $\zeta(t)$ der klassischen Bewegungsgleichung (\ref{dgl_zeta}) zu \cite{mads}
      \begin{equation}
        \zeta(t) = \frac{A\sin(\omega t)}{m(\omega_0^2 - \omega^2)} \; .
      \end{equation}
      Das Berechnen des Wirkungsintegrals und anschließendes Identifizieren des linearen Anteils liefert die Quasienergien für die gegebene Kraft:
      \begin{equation}
        \epsilon_n  = \hbar \omega_0\left(n+\frac{1}{2}\right) - \frac{A}{4m(\omega_0^2-\omega^2)} \;.
      \end{equation}
      Wie zu erkennen streben die Quasienergien, und somit die mittlere Energie des Systems (\ref{mittleres_H}), gegen Unendlich, wenn sich die Treibfrequenz $\omega$ nahe der Eigenfrequenz des Oszillators $\omega_0 $ befindet.
      Da wir einen getriebenen Oszillator ohne Dämpfung betrachten, war dieses Ergebnis zu erwarten \cite{mads}.



























  b

\chapter{Ergebnisse}


\section{Quasienergien des einzelnen getriebenen Oszillators f"ur eine beliebige periodische Treibkraft $S(t)$}
  Um die Quasienergien einer beliebigen periodisch treibenden Kraft auf elegante Weise zu bestimmen, und ohne ein Wirkungsintegral berechnen zu m"ussen, setzen wir eine komplexe Furier-Reihe an:
  \begin{equation}
    S(t) = \sum_{j=-\infty}^\infty c_j \text e^{\text ij\omega t}, \; c_j = \frac{1}{T} \int_0^T S(t) \text e^{-ij\omega t} \: \text d t \; .
  \end{equation}
  Die $c_j$ haben die Einheit einer Kraft.

  F"ur $\zeta(t)$ w"ahlen wir jetzt ebenfalls einen Reihenansatz:
  \begin{equation}
    \zeta(t) = \sum_{j=-\infty}^\infty d_j \text e^{\text ij\omega t} \; .
  \end{equation}
  Wir betrachten wieder nur die inhomogene klassische Bewegungsgleichung und bestimmen die $d_j$, indem wir einsetzen und einen Koeffizientenvergleich der $d_j$ und $c_j$ machen.
  Es zeigt sich, dass
  \begin{equation}
    \zeta(t) = \sum_{j=-\infty}^\infty \frac{c_j}{m(-j^2\omega^2+\omega_0^2)} \text e^{ij\omega t}
  \end{equation}
  gilt, womit sich die Lagrange-Funktion (\ref{lagrange_zeta}) ergibt zu:
  \begin{align}
    \begin{split}
      L &= \frac{1}{2}m\dot \zeta^2 - \frac{1}{2}m\omega_0^2\zeta^2 + S(t)\zeta \\
       &=\sum_j \sum_l \left[ \frac{-\omega^2}{2m} \frac{jc_j}{-j^2\omega^2+\omega_0^2} \frac{lc_l}{-l^2\omega^2+\omega_0^2}
       -\frac{\omega_0^2}{2m} \frac{c_j}{-j^2\omega^2+\omega_0^2} \frac{c_l}{-l^2\omega^2+\omega_0^2} \right. \\
        &\left. \quad + \frac{1}{m} \frac{c_jc_l}{-j^2\omega^2+\omega_0^2} \right] \text e^{\text i(j+l)\omega t}  \; , \; j,l \in \mathbb{Z} \; .
     \end{split}
   \end{align}
  \iffalse
  \begin{align}
    \begin{split}
      L &= \frac{1}{2}m\dot \zeta^2 - \frac{1}{2}m\omega_0^2\zeta^2 + S(t)\zeta \\
       &=\frac{-\omega^2}{2m} \sum_j \sum_l \frac{jc_j}{-j^2\omega^2+\omega_0^2} \frac{lc_l}{-l^2\omega^2+\omega_0^2} \text e^{\text i(j+l)\omega t}\\
       &\quad-\frac{\omega_0^2}{2m} \sum_j \sum_l \frac{c_j}{-j^2\omega^2+\omega_0^2} \frac{c_l}{-l^2\omega^2+\omega_0^2} \text e^{\text i(j+l)\omega t}\\
       &\quad + \frac{1}{m} \sum_j \sum_l \frac{c_jc_l}{-j^2\omega^2+\omega_0^2} \text e^{\text i(j+l)\omega t}\; , \; j,l \in \mathbb{Z} \; .
     \end{split}
   \end{align}
   \fi

   Nun identifizieren wir alle Terme der Lagrange-Funktion, welche nach der Ausf"uhrung des Wirkungsintegrals linear in der Zeit $t$ sind, ohne dieses explizit zu berechnen.
   Denn wir wissen, dass die Exponentialterme, welche sich beim Integrieren nach der Zeit reproduzieren, nicht linear in $t$ sein k"onnen.
   Daher werden nur die konstanten Terme, welche entstehen wenn die Exponentialterme wegfallen, eine lineare Abh"angigkeit aufweisen.
   Die Exponentialterme werden 1 wenn $j=-l$ ist.
   Infolgedessen wird die Doppelsumme, die bei der Quadrierung entstanden ist, wieder zur Einzelsumme.
   Wir fassen den linearen Teil des Wirkungsintegrals zusammen:
   \begin{align}
     \begin{split}
       &\frac{1}{T} \int_0^T L \: \text d t \\
       &= \sum_j \left[\frac{\omega^2}{2m} \frac{j^2c_jc_{-j}}{(-j^2\omega^2+\omega_0^2)^2}
       - \frac{\omega_o^2}{2m} \frac{c_jc_{-j}}{(-j^2\omega^2+\omega_0^2)^2} \right.
       \left. +\frac{1}{m} \frac{c_jc_{-j}}{-j^2\omega^2+\omega_0^2}
       \right]\\
       &= \sum_j \left[\frac{1}{2m} \frac{c_jc_{-j}(j^2\omega^2-\omega_0^2)}{(-j^2\omega^2+\omega_0^2)^2} + \frac{1}{m} \frac{c_jc_{-j}}{-j^2\omega^2+\omega_0^2} \right] \\
       &= \sum_j \frac{c_jc_{-j}}{2m(-j^2\omega^2+\omega_0^2)} = \frac{c_0^2}{2m\omega_0^2} + \sum_{j=1}^\infty \frac{c_j^2}{m(-j^2\omega^2+\omega_0^2)}
     \end{split}
   \end{align}

   Die Quasienergien $\epsilon_n$ sind hiernach:
   \begin{align}
     \begin{split}
       \epsilon_n &= \hbar \omega_0(n+\frac{1}{2}) - \sum_j \frac{c_jc_{-j}}{2m(-j^2\omega^2+\omega_0^2)}
     \end{split}
   \end{align}
   Interessanterweise kommt es in obiger Formel nicht nur bei $\omega = \omega_0$ zu einer Singularit"at, wie bei der Beispielkraft $S(t) = A\sin(\omega t)$, sondern bei allen $\omega = \omega_0 / j$.



   \subsection{Beispiele f"ur S(t) dreieck rechteck delta}





\section{Zwei getriebene gekoppelte Oszillatoren}
  In diesem Teil der Arbeit wird mit Hilfe der aus (\ref{lsg_einzelner}) bekannten L"osung des einzelnen getrieben Oszillators, die Wellenfunktionen f"ur ein System hergeleitet, das aus zwei gekoppelten Oszillatoren $x_1$ und $x_2$ der gleichen Masse $m$ besteht, von denen einer mit einer periodischen Kraft $S(t) = S(t+T)$ angetrieben wird.
  Die Potentialkonstanten $k$ der beiden Oszillatoren sind ebenfalls identisch, die Koplungskonstante $\kappa$ zwischen den Oszillatoren ist allerdings anders.
  Der Hamilton-Operator dieses Systems kann direkt aus der klassichen Mechanik "ubernommen werden:
  \begin{equation}
    H(t) = H(t+T) = \frac{p_1^2}{2m} + \frac{p_2^2}{2m} + \frac 1 2 kx_1^2 + \frac 1 2 kx_2^2 + \frac 1 2 \kappa(x_2-x_1)^2 - S(t)x_1 \; .
  \end{equation}
  Es werden auch die Erwartungswerte f"ur  den Ort $\braket{x_{1,2}}_{n.l}$ und den Impuls $\braket{p_{1,2}}_{n.l}$, genauso wie der Erwartungswert der Energie $\braket{H}_{n,l}$ und dessen zeitliches Mittel $\bar{H}_{n,l}$ berechnet und visualisiert.

  Zur L"osung des Systems wird eine unit"are Koordinatentransformation eingef"uhrt, welche den Hamilton-Operator $H(x_1,x_2,p_1,p_2,t)$ zu zwei in den neuen Koordinaten unabh"h"angigen Hamilton-Operatoren $H_+(x_+,p_+,t)$ und $H_-(x_-,p_-,t)$ mit effektiven Potentialkonstanten $k_+,k_-$ entkoppelt.
  Dann ergeben sich die Wellenfunktionen leicht aus denen des einzelnen getriebenen Oszillators.



  \subsection{Schr"odinger-Gleichung mit unabh"angigen Hamilton-Operatoren}
    Liegt ein Hamilton-Operator der Form
    \begin{equation}
      H = \sum_i H_i(x_i,p_i,t) \;,\; H_i : \cal H_i \rightarrow H_i
    \end{equation}
    vor, f"uhrt der Ansatz
    \begin{equation}
      \Psi = \prod_i \Psi(x_i,p_i,t) \; , \; \Psi_i \in \cal H_i
    \end{equation}
    auf unabh"angige Schr"odinger-Gleichungen f"ur die einzelnen Wellenfunktionen $\Psi_i(x_i,p_i,t)$ \cite{online quelle}.
    Um dies schnell zu zeigen, schauen wir uns den Fall von zwei unabh"angigen Hamiltonions an:
    \begin{align}
      \begin{split}
        \text i \hbar \frac{\partial}{\partial t} \Psi = H\Psi \iff \Psi_2 H_1 \Psi_1 + \Psi_1 H_2 \Psi_2 = \text i \hbar \left(\frac{\partial}{\partial t} \Psi_1\Psi_2 + \Psi_1\frac{\partial}{\partial t} \Psi_2 \right) \; .
      \end{split}
    \end{align}
    Da die Operatoren nur auf Funktionen wirken, die auf dem selben Raum operieren, kann man sie an der jewails anderen Funktion vorbei ziehen.
    Wenn wir weiterhin durch den $\Psi_1\Psi_2$ teilen, wird die Gleichung zu:
    \begin{equation}
      \frac{1}{\Psi_1}H_1\Psi_2 + \frac{1}{\Psi_2}H_2\Psi_1 = \text i \hbar \left(\frac{\frac{\partial}{\partial t} \Psi_1}{\Psi_1} + \frac{\frac{\partial}{\partial t} \Psi_2}{\Psi_2} \right) \; .
    \end{equation}
    Weil die linke und rechte Seite fur alle unabh"angigen $x_1,p_1,x_2,p_2$ gleich sein m"ussen, folgen die einzelnen Schr"odinger-Gleichungen:
    \begin{equation}
      \frac{1}{\Psi_1}H_1\Psi_2 = \frac{\frac{\partial}{\partial t} \Psi_1}{\Psi_1} \; , \; \frac{1}{\Psi_2}H_2\Psi_11 = \frac{\frac{\partial}{\partial t} \Psi_2}{\Psi_2} \; .
    \end{equation}
































   b

\chapter{Zusammenfassung und Ausblick}


\appendix
% Hier beginnt der Anhang, nummeriert in lateinischen Buchstaben
\chapter{Ein Anhangskapitel}

Hier könnte ein Anhang stehen, falls Sie z.B. Code, Konstruktionszeichnungen oder Ähnliches mit in die Arbeit bringen wollen. Im Normalfall stehen jedoch alle Ihre Resultate im Hauptteil der Bachelorarbeit und ein Anhang ist überflüssig.


\backmatter
\printbibliography

\cleardoublepage
\thispagestyle{empty}
\section*{Eidesstattliche Versicherung}
Ich versichere hiermit an Eides statt, dass ich die vorliegende Abschlussarbeit mit dem Titel \enquote{\thetitle} selbstständig und ohne unzulässige fremde Hilfe erbracht habe.
Ich habe keine anderen als die angegebenen Quellen und Hilfsmittel benutzt, sowie wörtliche und sinngemäße Zitate kenntlich gemacht. 
Die Arbeit hat in gleicher oder ähnlicher Form noch keiner Prüfungsbehörde vorgelegen.

\vspace*{1cm}\noindent
\begin{center}
  \begin{tabular}{@{}p{0.4\textwidth}@{\hspace{0.15\textwidth}}p{0.4\textwidth}@{}}
  \rule{\linewidth}{0.25pt}& \rule{\linewidth}{0.25pt}\\
  Ort, Datum & Unterschrift
  \end{tabular}
\end{center}

\subsection*{Belehrung}
Wer vorsätzlich gegen eine die Täuschung über Prüfungsleistungen betreffende Regelung einer Hochschulprüfungsordnung verstößt, handelt ordnungswidrig.
Die Ordnungswidrigkeit kann mit einer Geldbuße von bis zu \SI[round-mode=places, round-precision=2]{50000}{€} geahndet werden. 
Zuständige Verwaltungsbehörde für die Verfolgung und Ahndung von Ordnungswidrigkeiten ist der Kanzler/die Kanzlerin der Technischen Universität Dortmund. 
Im Falle eines mehrfachen oder sonstigen schwerwiegenden Täuschungsversuches kann der Prüfling zudem exmatrikuliert werden \mbox{(\S\,63 Abs. 5 Hochschulgesetz --HG--).}

Die Abgabe einer falschen Versicherung an Eides statt wird mit Freiheitsstrafe bis zu 3 Jahren oder mit Geldstrafe bestraft.

Die Technische Universität Dortmund wird ggf.\ elektronische Vergleichswerkzeuge (wie z.\,B.\ die Software \enquote{turnitin}) zur Überprüfung von Ordnungswidrigkeiten in Prüfungsverfahren nutzen. \\[\baselineskip]

\noindent Die oben stehende Belehrung habe ich zur Kenntnis genommen.\\[1cm]
\begin{center}
\begin{tabular}{@{}p{0.4\textwidth}@{\hspace{0.15\textwidth}}p{0.4\textwidth}@{}}
\rule{\linewidth}{0.25pt}& \rule{\linewidth}{0.25pt}\\
Ort, Datum & Unterschrift
\end{tabular}
\end{center}

\end{document}
